\section{Convex Functions}

\begin{definition}[Convex Set]
    \label{def:convex-set}%
    A set \(S \subseteq \RR^n\) is \emph{convex}, if for any \(\vb{x}, \vb{y} \in S\) and \(t \in \ccintv{0}{1}\), we have
    \[
        t \vb{x} + \pqty{1 - t} \vb{y} \in S.
    \]
\end{definition}

\autoref{def:convex-set} says that a set is convex, if and only if the straight line segment that connects any two points are included in the set.

\begin{definition}[Graph of Function]
    \label{def:graph-of-fn}%
    Let \(f\) be a function with domain \(\dom\pqty{f} \subseteq \RR^n\), i.e., \(\functiondomain{f}{\dom\pqty{f}}{\RR}\). The \emph{graph} of \(f\) is the surface
    \[
        z = f\pqty{\vb{x}}
    \]
    in \(\RR^{n + 1}\), with coordinates \(\pqty{\vb{x}, z}\). A \emph{chord} of \(f\) is a line segment joining \(2\) points of its graph.
\end{definition}

\begin{definition}[Convex Function]
    \label{def:convex-function}%
    A function \(f\) is \emph{convex}, if
    \begin{itemize}
        \item its domain \(\dom\pqty{f}\) is convex (as a set); and
        \item the graph of \(f\) lies below, or on, its chords; mathematically, for any \(\vb{x}, \vb{y} \in S\) and \(t \in \ccintv{0}{1}\), we have
        \begin{equation}
            f\pqty{t \vb{x} + \pqty{1 - t} \vb{y}} \leq t f\pqty{\vb{x}} + \pqty{1 - t} f\pqty{\vb{y}}. \label{eq:convex-condition}
        \end{equation}
    \end{itemize}

    If instead, \autoref{eq:convex-condition} is strict: for any \(\vb{x}, \vb{y} \in S\) with \(\vb{x} \neq \vb{y}\), and \(t \in \oointv{0}{1}\), we have
    \begin{equation}
        f\pqty{t \vb{x} + \pqty{1 - t} \vb{y}} < t f\pqty{\vb{x}} + \pqty{1 - t} f\pqty{\vb{y}}, \label{eq:strict-convex-condition}
    \end{equation}
    then we say \(f\) is \emph{strictly convex}.

    We say \(f\) is \emph{(strictly) concave} if \(-f\) is (strictly) convex.
\end{definition}

Note that the first property is important, since we need the domain to be convex for \(f\pqty{t \vb{x} + \pqty{1 - t} \vb{y}}\) to be well-defined in the first place.

\begin{figure}[ht!]
    \centering
    \begin{tikzpicture}
        \begin{axis}[
            axis lines = middle,
            xlabel = {\(x\)},
            ylabel = {\(f\pqty{x}\)},
            xmin = 0, xmax = 6,
            ymin = 0, ymax = 6,
            samples = 100,
            domain = 0.5:5.5,
            xtick = {2, 3, 5},
            xticklabels = {\(x\), \(t x+\pqty{1-t}y\), \(y\)},
            ytick = \empty
        ]

            \addplot[no marks] {0.5*(x-3)^2 + 1};
            \node[above left] at (axis cs:5.5,4.125) {\(f\)};

            \addplot[only marks, mark=*, mark size=1pt]
                coordinates {(2, 1.5) (5, 3)};

            \node[below left] at (axis cs:2, 1.5) {\(\pqty{x, f\pqty{x}}\)};
            \node[above left] at (axis cs:5, 3) {\(\pqty{y, f\pqty{y}}\)};

            \addplot[dashed] coordinates {(2, 1.5) (5, 3)};
            \addplot[only marks, mark=*, mark size=1pt]
                coordinates {(3, 2) (3, 1)};
            \node[above left] at (axis cs:3, 2) {\(t f\pqty{x} + \pqty{1 - t} f\pqty{y}\)};
            \node[below right] at (axis cs:3, 1) {\(f\pqty{t x + \pqty{1 - t} y}\)};

            \addplot[dotted] coordinates {(3, 0) (3, 2)};

        \end{axis}
    \end{tikzpicture}

    \caption{Convex Function}
    \label{fig:convex-fn}
\end{figure}

\begin{examples}[Convex Functions]
    \label{ex:convex-fn}%
    \begin{enumerate}
        \item The function \(f\pqty{x} = x^2\) is strictly convex, since
        \[
            \bqty{\pqty{1 - t} x + ty}^2 - \pqty{1 - t} x^2 - t y^2 = - t \pqty{1 - t} \pqty{x - y}^2 < 0
        \]
        for \(t \in \oointv{0}{1}\), \(x \neq y\).

        \item The functions \(f\pqty{x} = e^x\) and \(g\pqty{x} = e^{-x}\) are strictly convex.
        
        \item The function \(f\pqty{x} = \abs{x}\) is convex but not strictly convex: to show it is not strict, take for example \(x = 1\) and \(y = 2\), because some chords lie on the graph.
    \end{enumerate}
\end{examples}

\subsection{First Order Conditions}

\begin{theorem}[First Order Condition]
    \label{thm:first-order-condition}%
    If \(f\) is differentiable, then \autoref{eq:convex-condition} is equivalent to
    \begin{equation}
        \forall \vb{x}, \vb{y} \in \dom\pqty{f}: f\pqty{\vb{y}} \geq f\pqty{\vb{x}} + \pqty{\vb{y} - \vb{x}} \vdot \grad f\pqty{\vb{x}} \label{eq:first-order-condition}
    \end{equation}

    Geometrically, this means that the graph of \(f\) lies on or above any tangent plane, as shown in \autoref{fig:convex-fn-tangents} for the one-dimensional case.
\end{theorem}

\begin{figure}[ht!]
    \centering
    \begin{tikzpicture}
        \begin{axis}[
            axis lines = middle,
            xlabel = {\(x\)},
            ylabel = {\(f(x)\)},
            xmin = -1, xmax = 5,
            ymin = -2, ymax = 9,
            xtick = \empty,
            ytick = \empty,
            domain = 0:4,
            samples = 100,
            width = 10cm
        ]

            \addplot[no marks, thick, blue] {0.5*x^2};

            \addplot[only marks, mark=*, mark size = 1pt] coordinates {(1, 0.5) (2, 2) (3, 4.5)};

            \addplot[domain=0.25:4.5] {1 * (x - 1) + 0.5};
            \addplot[domain=0.4:4.2] {2 * (x - 2) + 2};
            \addplot[domain=0.5:4.1] {3 * (x - 3) + 4.5};
        \end{axis}
    \end{tikzpicture}

    \caption{Convex Function and its Tangents}
    \label{fig:convex-fn-tangents}
\end{figure}

\begin{corollary}[Convex Functions have Stationary Points as Global Minimums]
    \label{cor:stationary-point-global-minimum}%
    If \(f\) is convex and differentiable, then any stationary point is a global minimum.
\end{corollary}

\begin{proof}
    If \(\vb{a}\) is a stationary point, \(\grad f\pqty{\vb{a}} = \vb{0}\), so \autoref{eq:first-order-condition} implies that \(f\pqty{\vb{y}} \geq f\pqty{\vb{a}}\), for all \(\vb{y} \in \dom\pqty{f}\).
\end{proof}

Now we prove the first-order condition.

\begin{proof}[that \autoref{eq:convex-condition} implies \autoref{eq:first-order-condition}]
    Let
    \[
        h\pqty{t} = \pqty{1 - t} f\pqty{\vb{x}} + t f\pqty{\vb{y}} - f\pqty{\pqty{1 - t} \vb{x} + t \vb{y}}.
    \]

    Then note that
    \[
        h'\pqty{0} = - f\pqty{\vb{x}} + f\pqty{\vb{y}} - \pqty{\vb{y} - \vb{x}} \vdot \grad f\pqty{\vb{x}}
    \]
    which is the difference of two sides in \autoref{eq:first-order-condition}.

    But note that by definition,
    \[
        h'\pqty{0} = \lim_{t \to 0} \frac{h\pqty{t} - h\pqty{0}}{t} = \lim_{t \to 0} \frac{h\pqty{t}}{t}.
    \]

    But note that \autoref{eq:convex-condition} gives that \(h\pqty{t} \geq 0\) for all \(t \in \oointv{0}{1}\), and hence \(h'\pqty{0} \geq 0\), giving \autoref{eq:first-order-condition}.
 
\end{proof}

\begin{proof}[that \autoref{eq:first-order-condition} implies \autoref{eq:convex-condition}]
    In this direction, note that \autoref{eq:first-order-condition} implies that
    \begin{align*}
        f\pqty{\vb{x}} &\geq f\pqty{\vb{z}} + \pqty{\vb{x} - \vb{z}} \vdot \grad f\pqty{\vb{z}} \\
        f\pqty{\vb{y}} &\geq f\pqty{\vb{z}} + \pqty{\vb{y} - \vb{z}} \vdot \grad f\pqty{\vb{z}}
    \end{align*}
    for all \(\vb{x}, \vb{y}\) and \(\vb{z} \in \dom\pqty{f}\).

    Taking a linear combination, we have
    \[
        \pqty{1 - t} f\pqty{\vb{x}} + t f \pqty{\vb{y}} \geq f\pqty{z} + \bqty{\pqty{1 - t} \vb{x} + t \vb{y} - \vb{z}} \vdot \grad f\pqty{\vb{z}}.
    \]

    We now choose \(\vb{z} = \pqty{1 - t} \vb{x} + t \vb{y}\), which implies \autoref{eq:convex-condition}.
\end{proof}

\begin{proposition}[Alternative First Order Condition]
    An alternative way to write \autoref{eq:first-order-condition} is that
    \begin{equation}
        \pqty{\vb{y} - \vb{x}} \vdot \pqty{\grad f\pqty{\vb{y}} - \grad f\pqty{\vb{x}}} \geq 0
        \label{eq:first-order-condition-alternative}
    \end{equation}
    for all \(\vb{x}, \vb{y} \in \dom\pqty{f}\).
\end{proposition}

In particular, for \(n = 1\), \autoref{eq:first-order-condition-alternative} states that \(\pqty{y - x} \pqty{f'\pqty{y} - f'\pqty{x}} \geq 0\), which implies \(f'\pqty{y} \geq f'\pqty{x}\) if \(y \geq x\), i.e., \(f'\) is monotonically increasing.

\begin{proof}[that \autoref{eq:first-order-condition} implies \autoref{eq:first-order-condition-alternative}]
    From \autoref{eq:first-order-condition}, we have
    \begin{align*}
        f\pqty{\vb{y}} &\geq f\pqty{\vb{x}} + \pqty{\vb{y} - \vb{x}} \vdot \grad f\pqty{\vb{x}}, \\
        f\pqty{\vb{x}} &\geq f\pqty{\vb{y}} + \pqty{\vb{x} - \vb{y}} \vdot \grad f\pqty{\vb{y}}.
    \end{align*}

    Adding these give \autoref{eq:first-order-condition-alternative}.
\end{proof}

\begin{proof}[that \autoref{eq:first-order-condition-alternative} implies \autoref{eq:first-order-condition}]
    We let \(\vb{z} = \pqty{1 - t} \vb{x} + t \vb{y}\).

    Then we have
    \begin{align*}
        f\pqty{\vb{y}} - f\pqty{\vb{x}} &= \bqty{f\pqty{\vb{z}}}_{t = 0}^{t = 1} \\
        &= \int_{0}^{1} \dd{t} \dv{t} f\pqty{\vb{z}} \\
        &= \int_{0}^{1} \dd{t} \pqty{\vb{y} - \vb{x}} \vdot \grad f\pqty{\vb{z}}.
    \end{align*}

    Hence,
    \[
        f\pqty{\vb{y}} - f\pqty{\vb{x}} - \pqty{\vb{y} - \vb{x}} \vdot \grad f\pqty{\vb{x}} = \int_{0}^{1} \dd{t} \Bqty{\pqty{\vb{y} - \vb{x}} \vdot \bqty{\grad f\pqty{\vb{z}} - \grad f \pqty{\vb{x}}}}.
    \]

    Replace \(\vb{y}\) with \(\vb{z}\) in \autoref{eq:first-order-condition-alternative} gives that \(t \Bqty{\cdots} \geq 0\) (by using that \(\vb{z} = \pqty{1 - t} \vb{x} + t \vb{y}\)), and hence \(\Bqty{\cdots} \geq 0\).

    Hence, the integral is non-negative, and \autoref{eq:first-order-condition} holds.
\end{proof}

\subsection{Second Order Condition}

\begin{theorem}[Second Order Condition]
    \label{thm:second-order-condition}%
    If \(f\) is \(\contClass{2}\), then \autoref{eq:convex-condition} is equivalent to all eigenvalues of the Hessian \(\vb{H} \pqty{\vb{x}}\) is non-negative, for all \(\vb{x} \in \dom\pqty{f}\). In other words, the Hessian is positive semi-definite.
\end{theorem}

\begin{proof}[that \autoref{eq:convex-condition} implies \autoref{thm:second-order-condition}]
    We start from \autoref{eq:first-order-condition-alternative}. Set \(\vb{Y} = \vb{x} + \vb{h}\) in \autoref{eq:first-order-condition-alternative}, we have
    \[
        \vb{h} \vdot \bqty{\grad f\pqty{\vb{x} + \vb{h}} - \grad f\pqty{\vb{x}}} \geq 0.
    \]

    By Taylor expanding, we have
    \[
        \nabla_i f\pqty{\vb{x} + \vb{h}} = \nabla_i f \pqty{\vb{x}} + h_j H_{ij} \pqty{\vb{x}} + \bigO \pqty{h^3} \geq 0
    \]
    and hence
    \[
        h_i h_j H_{ij} \pqty{\vb{x}} + \bigO \pqty{h^3} \geq 0.
    \]

    Suppose that \(\vb{H}\pqty{\vb{x}}\) has some negative eigenvalue \(\lambda\) with eigenvector \(\vb{e}\), then setting \(\vb{h} = h \vb{e}\) gives that
    \[
        \lambda h^2 \vb{e}^2 + \bigO \pqty{h^3} \geq 0
    \]
    but the left-hand side is negative for sufficiently small \(h\) (since \(\lambda < 0\)), which is a contradiction.
\end{proof}

\begin{proof}[ththat \autoref{thm:second-order-condition} implies \autoref{eq:convex-condition}]
    We only prove in the case \(n = 1\) in this case, and we have \(H\pqty{x} = f''\pqty{x}\) and hence \(f''\pqty{x} \geq 0\) for all x.

    We have
    \[
        0 \leq \sign \pqty{y - x} \int_{x}^{y} f'' \pqty{x} \dd{x} = \sign \pqty{y - x} \pqty{f'\pqty{y} - f'\pqty{x}}
    \]
    and hence multiply by \(\abs{y - x}\), we have
    \[
        0 \leq \pqty{y - x} \pqty{f'\pqty{y} - f'\pqty{x}}
    \]
    which gives the result for \autoref{eq:first-order-condition-alternative}, hence \autoref{eq:first-order-condition}.
\end{proof}

\begin{example}[Using Second Order Conditions]
    Consider the function \(f\pqty{x, y} = \flatfrac{1}{xy}\), and \(\dom\pqty{f} = \set{\pqty{x, y}}{x > 0, y > 0}\). The Hessian is given by
    \[
        \vb{H} = \frac{1}{xy} \pmqty{\frac{2}{x^2} & \frac{1}{xy} \\ \frac{1}{xy} & \frac{2}{y^2}},
    \]
    and we have
    \[
        \det \vb{H} = \frac{3}{x^4 y^4} > 0
    \]
    and
    \[
        \tr \vb{H} = \frac{2}{xy} \pqty{\frac{1}{x^2} + \frac{1}{y^2}} > 0.
    \]

    This means the eigenvalues of \(\vb{H}\) is strictly true in \(\dom\pqty{f}\), and hence \(f\) is strictly convex.

    Note that if we take the domain to be
    \[
        \dom\pqty{f} = \set{\pqty{x, y}}{xy > 0},
    \]
    it seems like \(f\) would still be convex -- but that is not the case, because \(D\pqty{f}\) is not convex, and hence \(f\) is not convex,
\end{example}